\section{Polynomial interpolation, regularity, and multiplicity}
\label{poly:section}

We record the algebraic facts used below, with their degree dependence.  The
ambient dimension in this section is denoted by $n$ and is fixed.  Constants
written $C_n$ depend only on $n$; they may be increased between statements.
Varieties are reduced unless a local quotient or an intersection scheme is
explicitly specified.  The degree of an affine variety is the degree of its
projective closure.  For a family of distinct irreducible components of one
dimension, its degree is the sum of their degrees.  Components are counted
only once, even if they occur in more than one defining system.

The section supplies interpolation from defining equations, restriction-rank
bounds from component degrees, common equations with controlled exceptional
degree, and tests for high multiplicity and projection loss.  We first
prove isolated-point interpolation.  Applying it on generic fibers will
then give the rank and regularity estimates.

For a subset $T\subset\C^n$, put
\[
 H_T(D)=\dim_{\C}\bigl(\C[x_1,\ldots,x_n]_{\le D}/
           I(T)_{\le D}\bigr),\qquad D\in\mathbb Z_{\ge0}.
\]
Here $I(T)$ consists of the polynomials vanishing on $T$.  The degree-$D$
envelope is
\[
 \Env_D(T)=\{x:f(x)=0\text{ for every }f\in I(T)_{\le D}\}.
\]
Its defining collection can be replaced by a finite basis of
$I(T)_{\le D}$.  Directly from the definition,
\begin{equation}\label{poly:envelope-equality}
 H_{\Env_D(T)}(D)=H_T(D).
\end{equation}
For a projective set we use restrictions of homogeneous polynomials of
exact degree $D$.  Dehomogenizing in any affine chart meeting every
component gives the same rank: homogenization identifies polynomials of
degree at most $D$ with homogeneous polynomials of degree $D$, and a
polynomial vanishes on a component if and only if it vanishes on a dense
affine part of that component.  We always choose charts generically for
the finitely many points and components under discussion.

We use the following forms of the elementary intersection facts.  A
proper hypersurface section of a pure $k$-dimensional variety of degree
$a$, by a polynomial of degree $q$, has reduced support of pure dimension
$k-1$ and degree at most $aq$; the section may be empty in an affine
chart.  Successive proper sections obey the corresponding product bound.
More generally, the degrees, with local intersection multiplicities, of
isolated components of an intersection are at most the B\'ezout product,
even if the same equations have larger components elsewhere.  One obtains
this latter form by taking generic linear combinations locally at the
isolated components, perturbing to proper intersections, and specializing:
the isolated intersection multiplicities are positive and persist in the
specialization, whereas intersections on other components can only add
to the total.  We will never replace an isolated component by a
subvariety lying inside a larger component.

These intersection bounds are the componentwise refined B\'ezout
inequality \cite[Theorem~12.3 and Example~12.3.1]{fulton1998}.
For isolated complete intersections of hypersurfaces in a regular
ambient space, intersection multiplicity equals local quotient length.
The hypersurface factors are Cohen--Macaulay, as required in
\cite[Example~12.3.7(ii)--(iii)]{fulton1998}.

\subsection{Separators supplied by the distance set}

\begin{lemma}[Physical separators and sample bounds]\label{poly:separators}
Let $P\subset\R^n$ be finite and let $M=1+|\Delta(P)|$.  Each $p\in P$
has a polynomial $s_p$ of degree at most $2(M-1)$ such that
\[
 s_p(p)=1,\qquad s_p(q)=0\quad(q\in P\setminus\{p\}).
\]
If $V\subset\C^n$ is a pure $u$-dimensional variety of degree $a$, then
\begin{equation}\label{poly:physical-count}
 |P\cap V|\le C_n aM^u.
\end{equation}
Moreover, $P$ is precisely the common zero set of a system of polynomials
of degree at most $2M+1$.  In particular, all its points are isolated
components of that system.
\end{lemma}

\begin{proof}
For $p\in P$ define
\[
 s_p(x)=\prod_{\rho\in\Delta(P)}
       \frac{\sum_{j=1}^n(x_j-p_j)^2-\rho^2}{-\rho^2}.
\]
Every denominator is nonzero.  The stated values on the real sample
follow even though we subsequently regard $s_p$ as a polynomial over
$\C$.

We prove \eqref{poly:physical-count} by induction on $u$ and add over the
irreducible components of $V$.  For $u=0$ the number of points is at most
$a$.  Suppose $V$ is irreducible of positive dimension and contains a
sample point $p$.  The polynomial $s_p$ does not vanish identically on
$V$, because it is nonzero at $p$, and it vanishes at every other sample
point of $V$.  Its proper section has degree at most $2Ma$.  The induction
hypothesis therefore gives
\[
 |P\cap V|\le 1+C_n(2Ma)M^{u-1}.
\]
Increasing the dimension-dependent constant proves the claim.  This
argument requires no smoothness or general-position assumption on the
sample points.

Finally use the polynomials
\begin{equation}\label{poly:exact-isolation}
 1-\sum_{p\in P}s_p(x),\qquad
 (x_j-p_j)s_p(x)\quad(p\in P,\ 1\le j\le n).
\end{equation}
They vanish at every point of $P$.  At any common zero, the first equation
forces some $s_p$ to be nonzero, and the corresponding $n$ coordinate
equations then force $x=p$.  Thus there are no additional zeros.
\end{proof}

\subsection{Interpolation at isolated points}

The next lemma concerns the isolated points of a system, not a system
whose entire zero set is necessarily finite.  This distinction will be
needed for envelopes and for intermediate cuts.

\begin{lemma}[A determinant in a local parameter quotient]
\label{poly:socle}
Let $F$ be an algebraically closed field, let
$R=F[Z_1,\ldots,Z_c]_{(Z_1,\ldots,Z_c)}$, and suppose
$f_1,\ldots,f_c\in R$ generate an ideal of finite colength.  If
\[
 (f_1,\ldots,f_c)^{\mathsf t}
       =Q(Z)(Z_1,\ldots,Z_c)^{\mathsf t},
\]
then the class of $\det Q$ in $R/(f_1,\ldots,f_c)$ is nonzero.
\end{lemma}

\begin{proof}
The coordinate parameters form a regular sequence.  The regular local
ring $R$ is Cohen--Macaulay, so the $c$ elements $f_i$, being a system of
parameters, also form a regular sequence
\cite[Tags~00OQ, 00NQ, and 02JN]{stacksproject}.  Their Koszul complexes
are therefore free resolutions of $R/(f)$ and $R/(Z)=F$, respectively
\cite[Tag~062F]{stacksproject}.
The matrix $Q$ induces a chain map from the first Koszul resolution to
the second, lifting the quotient map $R/(f)\longrightarrow F$; its map
in top degree is multiplication by $\det Q$.

For completeness, the vanishing needed to see that this top map is
nonzero is especially simple.  The coordinate Koszul resolution gives
\[
 \operatorname{Ext}^i_R(F,R)=0\quad(i<c),\qquad
 \operatorname{Ext}^c_R(F,R)=F.
\]
Every finite-length $R$-module has a composition series with factors
$F$.  The long exact sequence for $\operatorname{Ext}$ therefore gives
$\operatorname{Ext}^i_R(K,R)=0$ for every finite-length $K$ and $i<c$.
Apply this to the kernel $K$ of $R/(f)\longrightarrow F$.  The long exact
sequence supplies an injection
\[
 \operatorname{Ext}^c_R(F,R)\lhook\joinrel\longrightarrow
       \operatorname{Ext}^c_R(R/(f),R).
\]
Under the identifications provided by the two Koszul resolutions, this
map sends the class of $1$ to the class of $\det Q$, up to an irrelevant
sign.  Its image is nonzero.
\end{proof}

\begin{lemma}[Isolated-point interpolation, including first derivatives]
\label{poly:interpolation}
Let $F$ be algebraically closed, and let $S\subset F^c$ be any selection
of distinct isolated points of the common zero set of polynomials of
degree at most $L$, where $L\ge1$.  There are polynomials of degree at
most $c(L-1)$ separating the points of $S$.  In fact a separator for one
point can be chosen to vanish at every other common zero of the original
system, including its positive-dimensional components.

Arbitrary values and all first partial derivatives can be prescribed
simultaneously on $S$ by a polynomial of degree at most $C_cL$.
Furthermore, $S$ itself is the exact zero set of equations of degree at
most $C_cL$.
\end{lemma}

\begin{proof}
Fix an isolated point $z$ of the original system.  Its local defining
ideal is primary to the maximal ideal at $z$.  We may choose $c$ constant
linear combinations $f_1,\ldots,f_c$ of the given equations whose local
ideal has finite colength.  To see this, choose the combinations
successively, avoiding the finitely many minimal primes of the preceding
proper intersections other than the maximal ideal.  The original ideal
is not contained in any of those primes.  Over an infinite field the
forbidden coefficient vectors form a finite union of proper linear
subspaces, so each choice is possible.  Each $f_i$ still has degree at
most $L$.  The same avoidance can be imposed simultaneously at any
finite list of isolated points.  Thus a generic choice of $c$
combinations is a local parameter sequence at every point of that
list, although the argument below only needs this property at $z$.

Introduce two sets of variables $Y,Z$.  Telescoping one coordinate at a
time gives a divided-difference matrix $Q(Y,Z)$ such that
\begin{equation}\label{poly:divided-difference}
 f(Y)-f(Z)=Q(Y,Z)(Y-Z).
\end{equation}
Every entry has degree at most $L-1$ in $Y$, and hence
$\det Q(Y,Z)$ has degree at most $c(L-1)$ in $Y$.  Take the coefficients
in the local Artin algebra
\[
 A_z=F[Z]_{(Z-z)}/(f_1(Z),\ldots,f_c(Z)).
\]
If $y\ne z$ is any common zero of the original system, then $f(y)=0$ and
one coordinate of $y-Z$ is a unit in $A_z$.  Multiplying
$Q(y,Z)(y-Z)=0$ by the adjugate of $Q(y,Z)$ shows that
\[
 \det Q(y,Z)=0\quad\hbox{in }A_z.
\]
This works for every such $y$, whether or not it is isolated.

On setting $Y=z$ in \eqref{poly:divided-difference}, we obtain
$f(Z)=Q(z,Z)(Z-z)$.  Lemma~\ref{poly:socle}, after translation, says
that $\det Q(z,Z)$ is nonzero in $A_z$.  Choose an $F$-linear functional
$\lambda:A_z\longrightarrow F$ which is nonzero on this element.
Applying $\lambda$ coefficientwise to $\det Q(Y,Z)$ and normalizing
produces a polynomial $s_z$ of the required degree, with $s_z(z)=1$ and
$s_z(y)=0$ at every other common zero.

Here are explicit first-order interpolants, which also show that the
number of points does not enter the degree bound.  The polynomial
\[
 u_z=3s_z^2-2s_z^3
\]
has value $1$ and all first derivatives zero at $z$, and has value and
all first derivatives zero at the other selected points.  The
polynomial
\[
 v_{z,j}=(Y_j-z_j)s_z^2
\]
has zero value at every selected point, derivative $1$ in coordinate $j$
at $z$, and all other prescribed first derivatives zero.  Linear
combinations of the $u_z$ and $v_{z,j}$ give arbitrary first-order data.
Their degrees are at most
$\max\{3c(L-1),2c(L-1)+1\}\le C_cL$.

The exact-isolation Equations \eqref{poly:exact-isolation}, with these
normalized $s_z$ and the set $S$, now have degree at most $c(L-1)+1$ and
have precisely $S$ as their common zero set.  The empty selection is
isolated by the equation $1=0$.
\end{proof}

\begin{remark}\label{poly:no-thickness-assumption}
The local quotient $A_z$ need not be reduced.  The determinant in
Lemma~\ref{poly:socle} records its nonzero top class, and a linear
functional extracts a polynomial which separates the underlying point.
Thus Lemma~\ref{poly:interpolation} does not assume that the original
equations have independent differentials, and it makes no assertion
about the thickness of other components.
\end{remark}

\begin{lemma}[A finite union of isolated-point lists]
\label{poly:union-interpolation}
Suppose $S$ is the union of finitely many lists $S_\nu\subset F^c$,
where $S_\nu$ consists of isolated points of a system of degree at most
$L_\nu\ge1$.  Put $L_*=\sum_\nu L_\nu$.  Values and first derivatives
on the distinct points of $S$ can be interpolated at degree at most
$C_cL_*$.
\end{lemma}

\begin{proof}
For a point $z\in S$, treat each list separately.  If $z\in S_\nu$,
take its normalized separator in that list.  If $z\notin S_\nu$,
the exact-isolation part of Lemma~\ref{poly:interpolation} supplies an
equation for $S_\nu$ which does not vanish at $z$; normalize its value
there to $1$.  The product over the lists is $1$ at $z$ and zero at every
other point of their union.  Its degree is at most $C_cL_*$.  The
first-order construction in the preceding proof costs a further fixed
factor only.
\end{proof}

\subsection{Component Hilbert ranks}

We now relate polynomial restriction rank to total component degree.
For the lower bound, a generic projection reduces the family to finitely
many points over the function field of the base.  Interpolation on that
fiber supplies independent transverse polynomials; multiplying them by
base monomials produces the required rank.

The upper estimate below is Chardin's Hilbert-function
bound~\cite{chardin1989}. For components of a single equation system,
the shifted lower estimate follows from the isolated-component bound
of Chardin--Philippon
\cite[Corollary~3]{chardinphilippon1999}; see also their
erratum~\cite{chardinphilippon2002erratum}. We give the proof, including
the extension to several equation systems used here.

\begin{lemma}[Upper and lower component ranks]\label{poly:hilbert}
Let $V\subset\C^n$ be a finite union of distinct irreducible varieties
of pure dimension $k$ and total degree $a$.  For every $D\ge0$,
\begin{equation}\label{poly:hilbert-upper}
 H_V(D)\le a\binom{D+k}{k}.
\end{equation}
Suppose in addition that every member of this family is an irreducible
component of a common locus defined by equations of degree at most
$L\ge1$.  Other components of that locus may have any dimension.  There
is an integer $q\le C_nL$ such that, for $D\ge q$,
\begin{equation}\label{poly:hilbert-lower}
 H_V(D)\ge a\binom{D-q+k}{k}.
\end{equation}
In particular,
\begin{equation}\label{poly:hilbert-at-L}
 H_V(L)\ge c_n aL^k.
\end{equation}
The same assertions hold for projective varieties, with homogeneous
restriction rank.

More generally, the lower assertions hold with $L$ replaced by
$L_* =\sum_\nu L_\nu$ if the family is a union of lists, the members of
the $\nu$th list being irreducible components of a system of degree at
most $L_\nu$.  A component in several lists is counted once.
\end{lemma}

\begin{proof}
We first justify the upper bound with its precise leading coefficient.
On a projective closure, polynomial restrictions inject into the global
sections of $\mathcal O_V(D)$.  A general hyperplane avoids the
associated components and has reduced pure-dimensional intersection
with a reduced variety; this is the reduced hyperplane-section form of
Bertini's Theorem in characteristic zero.  To see the reduced-scheme
assertion, the universal hyperplane incidence is a projective-space
bundle over the reduced variety, hence is reduced.  Its generic fiber
over the dual projective space is reduced by localization and is
geometrically reduced in characteristic zero.  The openness of reduced
fibers \cite[Tag~0578]{stacksproject} gives the assertion for a general
hyperplane.  Avoiding the components gives pure codimension one, and
the section has degree $a$.  The restriction exact sequence gives
\[
 h^0(V,\mathcal O_V(D))\le
 h^0(V,\mathcal O_V(D-1))+
 h^0(V\cap H,\mathcal O_{V\cap H}(D)).
\]
For dimension zero, the last space has dimension equal to the number
of reduced points, namely the degree.  In every dimension,
$h^0(V,\mathcal O_V)\le a$: it is at most the number of irreducible
components, since a regular function on an irreducible projective
variety is constant.  Induction on dimension, followed by summation in
$D$, gives
\[
 h^0(V,\mathcal O_V(D))
 \le a\sum_{j=0}^D\binom{j+k-1}{k-1}
 =a\binom{D+k}{k}
\]
for $k\ge1$.  This proves \eqref{poly:hilbert-upper}; alternatively one
may add the bound for the irreducible components.  Affine
homogenization gives exactly the same restriction rank.

For the lower bound, choose generic linear coordinates
\[
 (t_1,\ldots,t_k,y_1,\ldots,y_{n-k})
\]
such that projection to the $t$ coordinates is finite of degree
$\deg V_i$ on every irreducible component $V_i$ of the family.  This is
the usual generic linear normalization: in the projective closure the
projection center is disjoint from $V_i$, and a general fiber is a
complementary linear section.  In characteristic zero the generic
fibers are reduced.  Put $K=\C(t_1,\ldots,t_k)$ and let $F$ be an
algebraic closure of $K$.  The geometric generic fiber of the selected
family consists of $a$ distinct points of $F^{n-k}$.

These points are isolated in the generic fiber of the \emph{entire}
given system.  Indeed, if $W$ is any different irreducible component
of that system, then $V_i\cap W$ is a proper subvariety of $V_i$.
Its projection to the $k$-dimensional base is not dominant.  Thus no
point of the generic fiber of $V_i$ lies on $W$.  This argument also
handles components $W$ of dimension larger than $k$; their generic
fibers may be positive-dimensional, but they do not contain the
selected sheets.

The equations on the generic fiber have degree at most $L$ in $y$.
Lemma~\ref{poly:interpolation} says that evaluation on its $a$ selected
points is surjective from $F[y]_{\le q}$ for some $q\le C_nL$.
Since ordinary monomials span this space, we can select $a$ monomials
$m_1(y),\ldots,m_a(y)$ of degree at most $q$ whose evaluation columns
are linearly independent over $F$.  For each $i$, multiply $m_i$ by
all base monomials $t^\alpha$ with $|\alpha|\le D-q$.  The resulting
global polynomials have total degree at most $D$ and are linearly
independent on $V$.  In fact a relation over $\C$ would give
\[
 \sum_i\left(\sum_\alpha c_{i,\alpha}t^\alpha\right)
                m_i\big|_{\text{generic fiber}}=0.
\]
Independence over $F$ forces each coefficient polynomial in $t$ to
vanish, and algebraic independence of the $t_j$ forces every
$c_{i,\alpha}$ to be zero.  Counting these monomials proves
\eqref{poly:hilbert-lower}.  Notice that no rational coefficient has
been cleared: only ordinary transverse monomials were selected.

We explain the passage to degree $L$, since the lower bound above
initially starts at a constant multiple of $L$.  In a monomial order
refining total degree, the standard monomials of $I(V)$ form a set
closed under division, and their number of degree at most $D$ is
$H_V(D)$.  Fix an integer $b\ge1$.  The map
\[
 x^\alpha\longmapsto x^{\lfloor\alpha/b\rfloor}
\]
sends standard monomials of degree at most $bL$ to standard monomials
of degree at most $L$.  Each image has at most $b^n$ preimages, by
the possible coordinatewise remainders.  Hence
\begin{equation}\label{poly:rank-rescaling}
 H_V(bL)\le b^n H_V(L).
\end{equation}
Take a fixed $b$ large enough that $bL-q\ge L$ and use
\eqref{poly:hilbert-lower}.  This proves \eqref{poly:hilbert-at-L}.

For the union-of-systems assertion, take a generic base projection
simultaneously for the finitely many systems and components.  On its
geometric generic fiber, Lemma~\ref{poly:union-interpolation} supplies
the same surjectivity at degree $C_nL_*$.  The monomial argument and
\eqref{poly:rank-rescaling} are unchanged.  This proof does not claim
that a component of one system stays an isolated component of the
union of all the original loci; it uses the selected finite lists
on the generic fiber instead.
\end{proof}

\begin{corollary}[Mixed-dimensional envelopes]\label{poly:envelope}
Let $Z$ be a locus of equations of degree at most $L\ge1$, and let
$a_k$ be the total degree of any chosen subfamily of its
$k$-dimensional irreducible components.  For their union $V$,
\[
 H_V(L)\ge c_n\sum_{k=0}^n a_kL^k.
\]
For an arbitrary reduced union with dimension-by-dimension degrees
$a_k$, without a hypothesis on its equations,
\[
 H_V(D)\le\sum_{k=0}^n a_k\binom{D+k}{k}.
\]
In particular these statements apply to the components of
$\Env_L(T)$.
\end{corollary}

\begin{proof}
For the lower bound, $H_V(L)$ is at least the rank of each pure
dimensional subfamily.  Lemma~\ref{poly:hilbert} bounds each of those
below by $c_na_kL^k$; taking a maximum loses at most the fixed factor
$n+1$.  For the upper bound, restriction to the entire union injects
into the direct sum of its restrictions to the pure-dimensional
subfamilies.  Apply \eqref{poly:hilbert-upper} to each.
\end{proof}

\subsection{Common regular equations and pooled proper cuts}

The rank bounds have two uses.  They give common equations with
independent differentials on most components, and they combine separate
proper cuts into one polynomial.  Both conclusions measure the exceptional
family by its degree, so that later mass bounds can control the points
lying on it.

Tidor, Yu, and Zakharov formulate regular approximate complete
intersections whose Jacobian recovers the tangent space generically
\cite[Theorem~4.1]{tidoryuzakharov2026}.  The next lemma has a different
family-level interface: its equations vanish on the whole component
family, and one degree cutoff controls an explicit exceptional-degree
loss.  We prove this interface below.

\begin{lemma}[Regularity for a whole component family]
\label{poly:regularity}
Let $V$ be a pure $k$-dimensional family of total degree $a$ as in the
lower-bound part of Lemma~\ref{poly:hilbert}, with equation bound $L$.
The union-of-systems version is allowed, with $L=\sum_\nu L_\nu$.
If $D\ge C_nL$, then, except for components of total degree at most
\begin{equation}\label{poly:regularity-loss}
 C_n aL/D,
\end{equation}
the equations of degree at most $D$ vanishing on the \emph{whole}
family have Jacobian rank $n-k$ generically on every component.

There is a finite set of these equations and a polynomial of degree
at most $C_nD$, obtained as a linear combination of their maximal
Jacobian minors, which is nonzero on each retained component.  At
points of $V$ where this polynomial is nonzero, their common zero
set is locally a smooth $k$-fold, equal locally to $V$.  In
particular distinct members of the family cannot meet there.
\end{lemma}

\begin{proof}
There is nothing to prove when $k=n$: the only possible nonempty
family is the full ambient space, with codimension zero.  Otherwise
choose the generic base and transverse coordinates used in
Lemma~\ref{poly:hilbert}, and keep the notation $K=\C(t)$ and
$F=\overline K$.  Let $b$ be the total degree of the components on
which the asserted generic Jacobian rank fails.  Their geometric
generic fibers contribute $b$ of the $a$ sheets.

We give the vector-space argument carefully.  Write
\[
 W_D=\C[t,y]_{\le D},\qquad
 E_D:W_D\longrightarrow F^a,
       \quad f\longmapsto (f(t,z_i))_{i=1}^a.
\]
These are a vector space and a linear map over $\C$, and
\[
 \ker_{\C}E_D=I(V)_{\le D},\qquad
 \rank_{\C}E_D=H_V(D).
\]
No linearity of this truncated kernel over $K$ is asserted.

At a failing sheet $z_i$, the common kernel of the ambient gradients
of the equations in $I(V)_{\le D}$ has dimension at least $k+1$ over
$F$.  Its intersection with the $(n-k)$-dimensional vertical space
therefore contains a nonzero vector $v_i$.  The tangent space of
the component projects isomorphically to the base at a generic
sheet, so this vertical vector is transverse to the component.
Define a first derivative functional
\[
 J_i(f)=v_i\mathbin{\cdot}\partial_yf(t,z_i).
\]
Its coefficients may be algebraic functions of $t$.  They are
coefficients of this functional, not functions to be differentiated.
By construction $J_i$ annihilates $\ker_{\C}E_D$.  If $J$ collects
these $b$ functionals, then
\begin{equation}\label{poly:jet-factorization}
 \rank_{\C}(E_D,J)=\rank_{\C}E_D=H_V(D).
\end{equation}

First-order interpolation on the geometric generic fiber gives
surjectivity onto these $a$ values and $b$ transverse derivatives
from $F[y]_{\le q}$, where $q\le C_nL$.  Select $a+b$ ordinary
$y$-monomials $m_i$ of degree at most $q$ with linearly independent
augmented columns $w_i\in F^{a+b}$.  Vertical differentiation gives
the identity
\[
 (E_D,J)(t^\alpha m_i)=t^\alpha w_i
       \qquad(|\alpha|\le D-q).
\]
All these input polynomials lie in $W_D$.  A relation over $\C$
among their columns becomes an $F$-linear relation among the $w_i$,
with coefficients that are polynomials in $t$.  Independence over
$F$, followed by algebraic independence of $t$ over $\C$, makes
every coefficient zero.  Thus
\begin{equation}\label{poly:augmented-lower}
 (a+b)\binom{D-q+k}{k}\le H_V(D)
                  \le a\binom{D+k}{k}.
\end{equation}
For $D\ge2q$, the ratio of the two binomial coefficients is at most
$1+C_nq/D$; this follows by multiplying the $k$ factors
$(D+i)/(D-q+i)$.  Rearranging \eqref{poly:augmented-lower} proves
\eqref{poly:regularity-loss}.  When $k=0$ the binomial coefficients
are both $1$ and the argument gives $b=0$.

Choose a basis $f_1,\ldots,f_s$ of $I(V)_{\le D}$.  All
$(n-k)$-row Jacobian minors have degree at most $(n-k)(D-1)$.
On each retained component at least one is nonzero.  A generic
constant linear combination $g$ of all these minors is consequently
nonzero on every retained component, since the excluded coefficient
vectors are a finite union of proper linear subspaces.  Wherever
$g\ne0$, some $n-k$ of the $f_i$ have independent differentials.
Their zero locus is a smooth $k$-fold locally by the complex
implicit function theorem.  The pure $k$-dimensional set $V$ is
contained in it and has the same local dimension, so its reduced
germ is that smooth germ.  All remaining $f_i$ vanish on this germ
as well.  Two distinct irreducible components cannot share this
germ: an open common germ is Zariski dense in either component and
would make them equal.
\end{proof}

We have obtained regular equations vanishing on the entire family.
We next combine different proper cuts on its components, controlling
those components on which the combined polynomial vanishes identically.

\begin{lemma}[Pooling different proper cuts]\label{poly:pooling}
Let $V=\bigcup_iV_i$ be a pure $k$-dimensional family of total
degree $a$ satisfying the lower-bound hypotheses of
Lemma~\ref{poly:hilbert}, with bound $L$.  On each $V_i$ let
$h_i$ be a polynomial of degree at most $L'\ge1$ which does not
vanish identically on $V_i$.  If
$D\ge C_n\max\{L,L'\}$, there is a polynomial $g$ of degree at
most $D$ which vanishes on every proper cut $V_i\cap\{h_i=0\}$
and is not identically zero on any member except possibly a
subfamily of total degree at most
\begin{equation}\label{poly:pooling-loss}
 C_n aL'/D.
\end{equation}
The union-of-systems version of the family is again allowed.
\end{lemma}

\begin{proof}
For $k=0$ all the cuts are empty, and take $g=1$.  Suppose $k\ge1$
and let $C$ be the union of their reduced supports.  B\'ezout gives
$\deg C\le aL'$ in dimension $k-1$, so
\[
 H_C(D)\le aL'\binom{D+k-1}{k-1}.
\]
Let $U_D=I(C)_{\le D}$, and call a component $V_i$ forced if every
polynomial in $U_D$ vanishes identically on it.  If their union has
degree $b$, the kernel inclusion
$I(C)_{\le D}\subset I(\bigcup_{i\text{ forced}}V_i)_{\le D}$
implies
\[
 H_{\bigcup_{i\text{ forced}}V_i}(D)\le H_C(D).
\]
The lower bound of Lemma~\ref{poly:hilbert} applies to this
subfamily with the same $L$.  For $D\ge C_nL$ it bounds the left
side below by $c_n bD^k$, proving \eqref{poly:pooling-loss}.

For each unforced component, the polynomials in $U_D$ which vanish
identically on it form a proper linear subspace of $U_D$.  Choose
$g$ outside the finite union of these subspaces.  On increasing
the fixed lower bound for $D/L'$, the degree estimate makes
$b<a$ for a nonempty family, so such a $g$ is a nonzero polynomial.
It has precisely the required properties.
\end{proof}

\subsection{Dividing a degree bound by a multiplicity}

We use multiplicity of a polynomial at a point in the elementary sense
of order of vanishing: all terms below that total degree vanish after
translation to the point.  The next argument supplies an actual
polynomial of smaller degree, rather than an estimate for the size of
the set of multiple zeros.  The following matrix lemma supplies such
vanishing orders from a rank deficit.  Together the two statements turn
an exceptional rank condition into a lower-degree equation that remains
nonzero at a specified test point.

\begin{lemma}[Multiplicity compression]\label{poly:compression}
Let $S\subset\C^n$ be finite and let $x_*\in\C^n$.  Suppose a polynomial
$f$ of degree at most $E$ has order at least $r\ge1$ at every point of
$S$, and $f(x_*)\ne0$.  There is a polynomial $g$ such that
\[
 \deg g\le (n+1)E/r,\qquad
 g|_S=0,\qquad g(x_*)\ne0.
\]
The degree bound has no dependence on $|S|$.
\end{lemma}

\begin{proof}
An empty set $S$ is handled by $g=1$.  Otherwise $E\ge r$, because
$f$ is a nonzero polynomial vanishing to order at least $r$ somewhere.
First work over an algebraically closed field of characteristic $p$,
where $p\ge r$.  Set
\[
 m=\left\lceil\frac{np}{r}\right\rceil.
\]
The Frobenius map is surjective on this field, so grouping monomials by
their exponent vectors modulo $p$ gives a unique expression
\begin{equation}\label{poly:frobenius-decomposition}
 f(X)^m=\sum_{0\le i_1,\ldots,i_n<p}X^i h_i(X)^p.
\end{equation}
Every nonzero $h_i$ has
\[
 \deg h_i\le mE/p\le nE/r+E/p\le(n+1)E/r.
\]

Fix $z\in S$.  The order of $f^m$ at $z$ is at least $mr\ge np$.
Every monomial in $X-z$ of total degree at least $np$ has some
coordinate exponent at least $p$.  It follows that
\[
 f^m\in J_z=((X_1-z_1)^p,\ldots,(X_n-z_n)^p).
\]
Modulo $J_z$, the identity $h_i(X)^p\equiv h_i(z)^p$ holds.  The
monomials $X^i$, with every digit below $p$, form a basis of the
quotient by $J_z$: translating the usual basis $(X-z)^i$ is an
invertible triangular change of basis.  Reducing
\eqref{poly:frobenius-decomposition} modulo $J_z$ therefore gives
$h_i(z)=0$ for every $i$.  This applies at every point of $S$.
Evaluation at $x_*$ shows that at least one $h_i(x_*)$ is nonzero,
because $f(x_*)^m\ne0$.  That $h_i$ proves the assertion in
characteristic $p$.

We include a concrete specialization argument to transfer this fixed
degree assertion to $\C$.  Put $q=\lfloor(n+1)E/r\rfloor$ and suppose
that no degree-$q$ polynomial separates the given $S=\{z_1,\ldots,z_s\}$
from $x_*$.  In the finite-dimensional space of polynomials of degree
at most $q$, this says that evaluation at $x_*$ belongs to the linear
span of the evaluations at the $z_j$.  Thus there are
$\lambda_1,\ldots,\lambda_s\in\C$ with
\begin{equation}\label{poly:failed-separation}
 x_*^\alpha=\sum_{j=1}^s\lambda_j z_j^\alpha
            \qquad(|\alpha|\le q).
\end{equation}
Let $R\subset\C$ be the finitely generated $\mathbb Z$-algebra generated
by the coefficients of $f$, all coordinates of these points, the
$\lambda_j$, and $f(x_*)^{-1}$.  All hypotheses are polynomial
identities over $R$: order at least $r$ is expressed by the vanishing
of the coefficients of total degree below $r$ in $f(z_j+Y)$.
The inverse in $R$ preserves $f(x_*)\ne0$ under every specialization.
The identities \eqref{poly:failed-separation} also hold in $R$.

Such an $R$ has specializations to algebraically closed fields of
arbitrarily large positive characteristic.  Here is the required
algebraic justification.  Noether normalization for the finitely
generated domain $R\otimes_{\mathbb Z}\mathbb Q$ gives algebraically
independent elements $t_1,\ldots,t_b$ over which it is finite.  After
inverting one nonzero integer $a$, these elements belong to $R[1/a]$
and the finitely many generators of $R[1/a]$ satisfy monic equations
over $\mathbb Z[1/a,t_1,\ldots,t_b]$.  Hence $R[1/a]$ is finite
integral over that polynomial ring.  For every prime $p\nmid a$, the
lying-over property supplies a maximal ideal above
$(p,t_1,\ldots,t_b)$.  Its residue field is a finite extension of
$\mathbb F_p$, and embeds in $\overline{\mathbb F}_p$.

Choose such a prime $p\ge r$.  Specializing the coefficients, points,
and identities gives characteristic-$p$ data satisfying the hypotheses
and \eqref{poly:failed-separation}.  The characteristic-$p$ construction
produces a degree-$q$ polynomial which vanishes at all the specialized
$z_j$ but not at the specialized $x_*$.  This contradicts
\eqref{poly:failed-separation}.  Repeated specialized points cause no
problem: the evaluation identity and the compression construction
remain valid.  The contradiction proves the complex assertion.
\end{proof}

\begin{lemma}[Rank loss gives ambient multiplicity]
\label{poly:matrix-multiplicity}
Let $A(x)$ be a matrix of polynomials on $\C^n$, and let $F(x)$ be a
$q$-row, $q$-column minor.  If $\rank A(p)\le q-h$, then
$\mult_pF\ge h$.  In particular, if $F$ is nonzero at a test point,
has degree at most $E$, and the indicated rank loss occurs at finitely
many points, Lemma~\ref{poly:compression} gives a polynomial of degree
at most $(n+1)E/h$ vanishing at those points but not at the test point.
\end{lemma}

\begin{proof}
Expand the determinant using columns from $A(p)$ and from
$A(x)-A(p)$.  A term with fewer than $h$ latter columns uses more
than $q-h$ constant columns and is zero, because those columns have
rank at most $q-h$.  Each remaining term belongs to the $h$th power
of the maximal ideal at $p$.  This is an assertion in the ambient
polynomial ring; no regularity of a subvariety containing the points
is needed.
\end{proof}

\subsection{Projection degrees and local multiplicities}

We next control the degree lost by projection from a point.  The exact
loss is the local multiplicity of the variety there.  Projection to a
hypersurface will relate this geometric multiplicity to polynomial
vanishing order, allowing the preceding compression lemma to be applied.

For an integral variety $X$ of dimension $k$ and $p\in X$, we write
$\mult_pX$ for its Hilbert--Samuel multiplicity.  Thus, with
$\mathfrak m_p$ the maximal ideal of its local ring,
\[
 \dim_{\C}\mathcal O_{X,p}/\mathfrak m_p^t
       =\frac{\mult_pX}{k!}\,t^k+O(t^{k-1}).
\]
Equivalently, this is the degree of its tangent cone, with the scheme
multiplicities of that cone included.  We set $\mult_pX=0$ for
$p\notin X$.  For a hypersurface this multiplicity is the order of
its local defining polynomial; for a smooth point it is $1$.

\begin{lemma}[Projection degree and multiplicity]\label{poly:projection}
Let $X\subset\PP^n_{\C}$ be irreducible of dimension $k\ge1$ and
degree $a$.
\begin{enumerate}[label=(\roman*)]
\item If projection from a point $p$ has image $Y$ of dimension $k$
and generic degree $t$, then
\begin{equation}\label{poly:point-projection-formula}
 t\deg Y=a-\mult_pX.
\end{equation}
\item If $k<n$, there is a generic linear projection to
$\PP^{k+1}$ which is finite and birational on $X$.  Its image is an
irreducible hypersurface of degree $a$.  Given finitely many specified
points $p_i\in X$, the projection can simultaneously be chosen so
that their fibers are singletons and
\begin{equation}\label{poly:projection-multiplicity}
 \mult_{\pi(p_i)}\pi(X)\ge\mult_{p_i}X.
\end{equation}
Given in addition a point $x_*\notin X$, it can also be required that
$\pi(x_*)\notin\pi(X)$.
\end{enumerate}
The assertions apply to affine varieties through their projective
closures and generic affine charts.
\end{lemma}

\begin{proof}
For (i), if $p\notin X$, projection is a morphism and the pullback of
a target hyperplane is a hyperplane section of $X$.  Intersecting $k$
general target hyperplanes and counting the $t$ points over a general
image point gives $t\deg Y=a$.

Suppose $p\in X$.  We give the intersection calculation that also
justifies the formula when $X$ is singular or not Cohen--Macaulay.
For the blowup construction and the degree formula for a linear system
with base locus, see \cite[Appendix~B.6 and Proposition~4.4]{fulton1998}.
Let $b:\widetilde X\longrightarrow X$ be the blowup of the ideal of
$p$, and let $E$ be its exceptional Cartier divisor.  The blowup is
the relative projective spectrum of the Rees algebra
$\bigoplus_{j\ge0}\mathcal I_p^j$, where $\mathcal I_p$ is the ideal
sheaf of $p$; it is the closure of the graph of
projection from $p$.  Write $H$ for the pullback of the hyperplane
divisor of $X$.  The linear forms through $p$ generate the ideal of
$p$ twisted by $\mathcal O_X(1)$, so the resolved projection
$f:\widetilde X\longrightarrow Y$ satisfies
\[
 f^*\mathcal O_Y(1)=\mathcal O_{\widetilde X}(H-E).
\]
Its top self-intersection is consequently $t\deg Y$, by taking
general hyperplanes on the target.

We recall each term in this self-intersection.  The divisor $H$
restricts trivially to $E$, since $E$ lies over the point $p$.
Therefore
\[
 H^{k-j}E^j=0\qquad(1\le j<k).
\]
Moreover
\[
 E=\operatorname{Proj}\!\left(
       \bigoplus_{j\ge0}\mathfrak m_p^j/\mathfrak m_p^{j+1}\right),
 \qquad \mathcal O(E)|_E=\mathcal O_E(-1).
\]
The Hilbert polynomial of this associated graded ring shows that
the degree of $E$ with respect to $\mathcal O_E(1)$ is
$\mult_pX$.  Thus $E^k=(-1)^{k-1}\mult_pX$.  Also $H^k=a$, as can
be computed using general hyperplanes avoiding $p$.  Expanding gives
\[
 t\deg Y=(H-E)^k=a+(-1)^kE^k=a-\mult_pX.
\]
Only Cartier-divisor intersection numbers on the integral blowup are
used here.  In particular, the calculation does not identify
multiplicity with the length of an arbitrary local linear section.

For (ii), the case $n=k+1$ is the identity projection.  Otherwise
choose a linear center $\Lambda\simeq\PP^{n-k-2}$.  For a fixed point
$q$, let $J(q,X)$ be the closure of the union of the lines joining
$q$ to points of $X$ different from $q$.  Its dimension is at most
$k+1$.  Choose one smooth point $q_0\in X$.  A general $\Lambda$
is disjoint from each of
\[
 X,\quad J(p_i,X),\quad J(q_0,X),\quad T_{q_0}X,
 \quad\text{and, if specified, }J(x_*,X).
\]
There are finitely many conditions, and each is a nonempty open
condition on the Grassmannian: the dimension of the relevant closed
set plus $n-k-2$ is strictly less than $n$.

The resulting projection defines a morphism $\pi:X\longrightarrow
\PP^{k+1}$.  It is finite.  Indeed, a fiber lies in a projective
linear space in which $\Lambda$ is a hyperplane; a positive-dimensional
projective fiber would meet that hyperplane, contradicting
$\Lambda\cap X=\varnothing$.  Avoidance of the joins makes the
fibers over the specified $p_i$, and over $q_0$, singletons.
Avoidance of $T_{q_0}X$ makes the differential injective at $q_0$,
so the fiber at $q_0$ is a reduced point.  More explicitly, its
Artinian local ring has zero cotangent space and hence zero maximal
ideal by Nakayama's Lemma.  The generic rank of a finite module is
at most the dimension of any fiber, again by Nakayama's Lemma.
The generic degree is therefore one.  The image is an integral
$k$-fold in $\PP^{k+1}$, hence a hypersurface, and its degree equals
$a$ by pulling back $k$ general hyperplanes.  The additional join
condition ensures $\pi(x_*)\notin\pi(X)$.

It remains to verify the local multiplicity assertion without a
smoothness assumption at $p_i$.  Write
\[
 A=\mathcal O_{\pi(X),\pi(p_i)},\qquad
 B=\mathcal O_{X,p_i},
\]
with maximal ideals $\mathfrak m$ and $\mathfrak n$.  Because the
fiber is a singleton, $B$ is a finite local $A$-algebra; because
the projection is birational, it has rank one and $A\subset B$.
The quotient $B/A$ is annihilated by a nonzero element of $A$ and
therefore has dimension at most $k-1$.  The Artin--Rees property
gives a fixed $c$ such that
\[
 \mathfrak m^jA\subseteq\mathfrak m^jB\cap A
                  \subseteq\mathfrak m^{j-c}A\qquad(j\ge c).
\]
Using the exact sequence
\[
 0\longrightarrow A/(\mathfrak m^jB\cap A)
 \longrightarrow B/\mathfrak m^jB
 \longrightarrow (B/A)/\mathfrak m^j(B/A)\longrightarrow0
\]
shows that $A/\mathfrak m^jA$ and $B/\mathfrak m^jB$ have the same
leading coefficient of degree $k$ in their lengths.  Since
$\mathfrak mB\subseteq\mathfrak n$,
\[
 \dim_{\C} B/\mathfrak m^jB\ge
       \dim_{\C}B/\mathfrak n^j.
\]
Comparison of the leading coefficients proves
\eqref{poly:projection-multiplicity}.
\end{proof}

\begin{corollary}[A bounded-degree test for high multiplicity]
\label{poly:high-multiplicity}
Let $X\subsetneq\C^n$ be irreducible of degree $a$, and let
$S\subset X$ be finite with $\mult_pX>a/2$ for every $p\in S$.
For any specified $x_*\notin X$, there is a polynomial of degree
at most $2(n+1)$ which vanishes on $S$ and is nonzero at $x_*$.
In particular the specified high-multiplicity centers lie in a
proper hypersurface of degree bounded in terms of the ambient
dimension alone.
\end{corollary}

\begin{proof}
For a zero-dimensional irreducible variety, $X$ is one reduced
point, and a linear polynomial suffices.  In positive dimension
apply Lemma~\ref{poly:projection}(ii) to the projective closure,
the finite set $S$, and $x_*$.  Compose a homogeneous defining
equation of the image with the linear forms of the projection and
dehomogenize in the original affine chart.  The resulting polynomial
has degree at most $a$ and is nonzero at $x_*$.  Choose the target
chart to contain the finitely many specified images.  Its
dehomogenization differs locally from this polynomial by a unit,
so its order at every $p\in S$ is at least $\mult_pX$: composition
does not decrease order, and hypersurface multiplicity equals the
order of its equation.  Apply Lemma~\ref{poly:compression} with
$r=\lfloor a/2\rfloor+1$.
\end{proof}

\begin{corollary}[A tangent hyperplane controls the projection loss]
\label{poly:tangent-projection}
Let an irreducible affine curve $C$ lie in a variety $X$ smooth at
$p\in C$.  Suppose $C$ has a point $z$ for which the direction
$z-p$ is not in $T_pX$.  Then
\[
 \mult_pC\le\tfrac12\deg C.
\]
Consequently, if projection of $C$ from $p$ is nonconstant, its
image degree times its generic map degree is at least
$\tfrac12\deg C$.
\end{corollary}

\begin{proof}
Choose an affine hyperplane through $p+T_pX$ which does not contain
$z$.  Its linear equation $\ell$ restricts to an element of
$\mathfrak m_{X,p}^2$, because its value and differential on $X$
both vanish at $p$.  On the curve it therefore belongs to
$\mathfrak m_{C,p}^2$, and it is not identically zero.

The one-dimensional reduced local ring of the curve is
Cohen--Macaulay.  Thus, writing $R=\mathcal O_{C,p}$, the local
intersection number with the hyperplane is
$\operatorname{length}(R/(\ell))$, and
$\operatorname{length}(R/(\ell^j))=
j\operatorname{length}(R/(\ell))$.
Since $(\ell^j)\subseteq\mathfrak m_{C,p}^{2j}$, comparison with
the Hilbert--Samuel leading term gives
\[
 \operatorname{length}(R/(\ell))\ge2\mult_pC.
\]
B\'ezout bounds this local intersection number by $\deg C$.
The final assertion follows from
\eqref{poly:point-projection-formula}.
\end{proof}

\begin{lemma}[The plane-curve polar bound]\label{poly:polar}
Let $C\subset\PP^2_{\C}$ be a reduced plane curve of degree $a$.
If $m_q=\mult_qC$, then
\begin{equation}\label{poly:polar-sum}
 \sum_{q:\,m_q\ge2}m_q(m_q-1)\le a(a-1),
 \qquad
 \sum_{q:\,m_q\ge2}m_q^2\le2a(a-1).
\end{equation}
The sums include points on any chosen line at infinity.
\end{lemma}

\begin{proof}
Write $C=\{F=0\}$ with $F$ homogeneous and squarefree.  A general
first polar
\[
 G=\lambda_0\frac{\partial F}{\partial X_0}
   +\lambda_1\frac{\partial F}{\partial X_1}
   +\lambda_2\frac{\partial F}{\partial X_2}
\]
has degree $a-1$ and has no irreducible factor in common with $F$.
Indeed, for every irreducible factor $F_i$ some derivative of $F_i$
is nonzero modulo $F_i$ in characteristic zero.  Modulo $F_i$, the
corresponding derivative of $F$ is this derivative times the product
of the other factors.  A generic coefficient vector avoids
vanishing identically on each of the finitely many components.

In an affine chart at $q$, write $f$ for the local equation of the
curve.  The three dehomogenized derivatives of $F$ are, after
renaming the coordinates,
\[
 f_x,\qquad f_y,\qquad af-xf_x-yf_y.
\]
Each has order at least $m_q-1$, so the local equation of $G$ has
order at least $m_q-1$.  The elementary local intersection
inequality gives
\[
 I_q(F,G)\ge m_q(m_q-1).
\]
One can see that inequality directly on the reduced local branches:
a branch of multiplicity $b$ has a primitive parametrization whose
two coordinates have minimum order $b$, and substituting into a
polynomial of order $c$ gives order at least $bc$.  Add over the
branches to obtain the displayed bound.  Since $F$ and $G$ have no
common component, B\'ezout gives total intersection number $a(a-1)$
and in particular makes the set of points with $m_q\ge2$ finite.
This proves the first inequality.  The second follows from
$m^2\le2m(m-1)$ for every integer $m\ge2$.
\end{proof}
